\documentclass[10pt,a4paper]{amsart}
\usepackage[margin=19mm]{geometry}
\usepackage{amsmath,amssymb,mathtools}
\usepackage[scr=rsfs,cal=euler]{mathalfa}
\usepackage{xfrac}
\usepackage[unicode,hidelinks,bookmarks=false]{hyperref}
\newcommand{\ie}{i.e.\ }
\newcommand{\cf}{cf.\ }
\newcommand{\resp}{respectively\ }
\newcommand{\Subsection}{\subsection}
\providecommand{\nobreakdash}{\nobreak\hbox{-}}
\setlength{\emergencystretch}{2em}
\multlinegap0pt
\usepackage{tikz}
\usetikzlibrary{cd}
\usepackage{mathtools}
\usepackage{dsfont}
\usepackage{amssymb}
\usepackage{xcolor}
\usepackage{tikz}
\newtheorem{lem}{Lemma}
\newtheorem{prop}{Proposition}
\newcommand{\C}{\mathcal{C}}
\DeclareMathOperator{\Conf}{Conf}
\DeclareMathOperator{\UConf}{UConf}
\DeclareMathOperator{\id}{id}
\title{The Barratt--Priddy--Quillen theorem via scanning methods}
\author{Marie-Camille Delarue}
\date{Source: arXiv:2510.13564v2 (CC BY 4.0)}
\begin{document}
\maketitle

\noindent\emph{Source and licence.} The Barratt--Priddy--Quillen theorem via scanning methods. Version arXiv:2510.13564v2 (CC BY 4.0), distributed by arXiv under the Creative Commons Attribution 4.0 International licence, http://creativecommons.org/licenses/by/4.0/, as recorded in the arXiv licence metadata for this exact version and shown on the abstract page https://arxiv.org/abs/2510.13564v2. Full licence notice: \texttt{licenses/arxiv-2510.13564-CC-BY-4.0.txt}.\par\smallskip

\noindent\textit{Editorial scope.} Counted unit: the article's prop deloopinglevelone (source0.tex lines 780--783). The article's complete original proof of it is delivered with it (source0.tex lines 785--910, one \texttt{proof} environment of 126 lines). No mathematics is written, repaired or re-derived: the statement and the proof are the article's own lines, in the article's own notation, and no retained line is substituted or re-flowed.  Retained uncounted as the source-local context the proof needs: lines 470--472; lines 736--739.  Nothing was omitted from any retained span, so the package carries no omission record.  No retained line contains a citation, so the wrapper carries no bibliography and no bibliography entry.  The retained lines contain the article's own diagram code, which the wrapper typesets with the standard tikz packages; no image file is delivered and the package declares no asset dependency.  Editorial formatting only, in addition: an AMS \texttt{amsart} preamble that restates the article's own declarations and macros used by the retained lines, namely the environments \texttt{lem}, \texttt{prop} on separate counters, together with the standard packages those lines need.  The retained environments print in this document as Lemma 1, Proposition 1 in order of appearance, whereas the article prints the same items with the numbering of its own full document.  The complete native source of the article as distributed in the arXiv e-print for this exact version is bundled unchanged as \texttt{sources/187123/source0.tex}.
\begin{lem}\label{goodnerve}
Let $N\in\mathbb{N}$ or $N=\infty$. The nerve of the category $\C_N$ is good (i.e., the category $\C_N$ is well-pointed).
\end{lem}

\begin{lem}\label{posetthing}
    Let $P$ be a subposet of $\mathbb{R}$.
    The classifying space of the topological poset $P$ is contractible.
\end{lem}

\begin{prop}\label{deloopinglevelone}
For $N\geq 1$, there is a zigzag of weak equivalences of $E_N$-algebras:
\[\Phi_0^N \simeq B\C_N.\]
\end{prop}

\begin{proof}
Let us construct a semi-simplicial resolution $X_\bullet$ of $\Phi_0^N.$
The space of $p$-simplices $X_p$ is given by $\Phi_0^N\times N_p\mathbb{R}$, where $\mathbb{R}$ is seen as a topological poset.
The face maps are given by, for $0\leq k \leq p$:
\begin{align*}
    d_k: X_p&\rightarrow X_{p-1}\\
    (\phi,t_0,\dots, t_p) &\mapsto (\phi,t_0,\dots, \hat{t_k},\dots, t_p).
\end{align*}
There is a semisimplicial map $\eta: X_\bullet\rightarrow N_\bullet\C_N$, which is defined on $p$-simplices as the map which sends $(\phi,t_0,\dots, t_p)$ to $\left((\phi_{\vert_{[t_0,t_1]}},t_0,t_1),\dots,(\phi_{\vert_{[t_{p-1},t_p]}},t_{p-1},t_p)\right)$. 
Here, $\phi_{\vert_{[t_{i-1},t_i]}}$ is an element of $\Phi_0^N([t_{i-1};t_i])$.
(Restricting $\phi$ to compact subintervals of $\mathbb{R}$ is continous with respect to the limit topology).
Informally, the map $\eta$ forgets what happens before $t_0$ and after $t_p$ and restricts $\phi$ to $p$ subintervals in between.

The projection $X_p\rightarrow N_p\mathbb{R}$ sending $(\phi, t_0,\dots, t_p)$ to $(t_0,\dots, t_p)$ is a fibration. 
Since the target of the fibration is contractible, the space $X_p$ is equivalent to the fiber over a fixed point $(t_0,\dots, t_p)$.
Given a $p$-tuple $(t_0,\dots,t_p)$, an element $\phi$ in the fiber is entirely determined by its restriction to all the intervals $(-\infty,t_0],\dots,[t_p,\infty)$ provided their endpoints coincide pairwise.
This means that the fiber $X'_p$ of the projection can be expressed as:
\begin{align}
X'_p =\Phi_0^N((-\infty,t_0])\times_{\UConf(n,I^N)}\dots\times_{\UConf(n,I^N)}&\Phi_0^N([t_k,t_{k+1}]) \label{fiberprime}\\
\times_{\UConf(n,I^N)}\dots\times_{\UConf(n,I^N)}\Phi_0^N([t_p,\infty)) \nonumber.
\end{align}

The map $N_p\C_N\rightarrow N_p\mathbb{R}$ is also a fibration, so let us consider the following diagram:
\begin{center}
    \begin{tikzcd}
        X'_p \arrow{r}\arrow{d}&Y'_p\arrow{d}\\
        X_p\arrow{r}{\eta_p}\arrow{d} &N_p\C_N\arrow{d}\\
        N_p\mathbb{R}\arrow{r}{=}&N_p\mathbb{R}.
    \end{tikzcd}
\end{center}
The restriction of $\eta_p$ to $X'_p$ is a map into the space
\[Y'_p:= \Phi_0^N([t_0,t_1])\times_{\UConf(n,I^N)}\dots\times_{\UConf(n,I^N)}\Phi_0^N([t_{p-1},t_{p}])\]
which forgets the two outer terms of the product in expression \ref{fiberprime}. 
We therefore need to prove that this operation does not change the homotopy type of the product.

The target maps from each $\Phi\left((-\infty,t_0]\times I^N, ([-R,R]\cap(-\infty,t_0])\times I^N\right)$ -- resp. source maps from $\Phi\left([t_p,\infty)\times I^N, ([-R,R]\cap[t_p,\infty))\times I^N\right)$ -- induce maps from the limit:
\begin{align*}
p_0: \Phi_0^N((-\infty,t_0])&\rightarrow \UConf(n,I^N)\\
p_p: \Phi_0^N([t_p,\infty))&\rightarrow \UConf(n,I^N)
\end{align*}
which are also fibrations, so the two outer pullbacks in (\ref{fiberprime}) are homotopy pullbacks.

Let us show that the two outer terms in this (homotopy) pullback are equivalent to $\UConf(n,I^N)$
and so that removing them does not change the homotopy type.
Let us consider the map $q_0:\UConf(n,I^N)\rightarrow \Phi_0^N((-\infty,t_0])$ induced by the maps:
$q_R:\UConf(n,I^N)\rightarrow \Psi((-\infty,t_0]\times I^N, [-R,R]\times I^N)$ which send a configuration $x$ to the constant path $\phi:t\mapsto (x,t)$.
Clearly, $p_0\circ q_0 = \id$. The following is a homotopy between $q_0\circ p_0$ and $\id$:
\begin{align*}
H: \Phi_0^N((-\infty,t_0])\times I &\rightarrow \Phi_0^N((-\infty,t_0])\\
(\phi,t_0,s) &\mapsto (\gamma_s,t_0),
\end{align*}
where $\gamma_1$ is the constant path at $\phi(t_0)$ and for $s<1$:
$$\gamma_s:t \mapsto \begin{cases}
&\phi(t-t_0-\frac{s}{1-s}),\text{ if }t\in(-\infty,t_0-\frac{s}{1-s}];\\ 
&\phi(t_0),\text{ if }t\in[t_0-\frac{s}{1-s},t_0].
\end{cases}$$ 
Therefore, the map $p_0$ is an equivalence.
Similarly, the map $p_p$ is an equivalence.

This proves that the lift of $\eta_p$ to fibers is an equivalence.
The map $\eta_p$ is a homotopy equivalence for all $p$, and so the map $\eta$ induces an equivalence on weak realizations:
\[\vert\vert X_\bullet\vert\vert \simeq \vert\vert N_\bullet\C_N\vert\vert,\]
but the thick realization of the nerve is equivalent to the ordinary geometric realization by lemma \ref{goodnerve}, therefore:
\[\vert\vert X_\bullet\vert\vert \simeq B\C_N.\]

Because $\vert\vert X_\bullet\vert\vert = \Phi_0^N\times \mathbb{B}\mathbb{R}$, there is a projection map 
\begin{align}
\epsilon: \vert\vert X_\bullet\vert\vert\rightarrow \Phi_0^N 
\end{align}
which forgets the $(t_k)$.
This map is a fibration.
% Indeed, given an outer square 
% \begin{center}
%     \begin{tikzcd}
%     D^k\times\{0\} \arrow[rr,"A"] \arrow[dd] \arrow[dr] && \vert\vert X_\bullet\vert\vert\arrow[dd,"\epsilon"]\\
%     &D^k\times[0,\delta] \arrow[dl]\arrow[ur,dashrightarrow,"H"]\\
%     D^k\times[0,1]\arrow[rr,"B"] &&\Phi_0^N,
%     \end{tikzcd}
% \end{center}
% where $A(d) = (\phi(d),t_0(d),\dots, t_p(d),x(d))$ and $B(d,s) = (\psi(d,s))$,
% we can find a lift $H$ by setting $H(s) = (\psi(d,s),t_0(d),\dots, t_p)$ for some small $\delta$. 
% The idea is that if $\phi$ moves just a little bit, it still won't hit the walls at the $t_k$.

The fiber over a point $\phi\in\Phi_0^N$ is the classifying space of the topological poset $\mathbb{R}$, which we show to be contractible in lemma \ref{posetthing}.
The map $\epsilon$ has contractible fibers so it is a homotopy equivalence, yielding the following zigzag:
\[\Phi_0^N\xleftarrow{\simeq}\vert\vert X_\bullet\vert\vert \xrightarrow{\simeq} B\C_N.\]
Additionally, these maps are compatible with the $E_N$-structures on $\Phi_0^N$ and $\C_N$.
In particular, when $N=\infty$, we get:
\[\Phi_0^\infty\xleftarrow{\simeq}\vert\vert X_\bullet\vert\vert \xrightarrow{\simeq} B\C_\infty. \qedhere\]

% For every $N$, there is a map 
% \begin{align*}
% \Phi_0^N &\rightarrow \Phi((-\infty,0]\times I^N)\times_{\UConf(n,I^\infty)\times\mathbb{R}} N_p\C_N\times_{\mathbb{R}\times \Conf(n,I^\infty)} \Phi([0,\infty) \times I^N)
% \end{align*}
% which sends an element $(\phi,t_0,\dots, t_p)$ to
% \[\left((\phi(t-t_0)\vert_{\mathbb{R}^-}),(\phi\vert_{[t_0,t_1]},t_0,t_1),\dots,(\phi\vert_{[t_{p-1},t_p]},t_{p-1},t_p),\phi(t-t_0)\vert_{\mathbb{R}^+}\right).
% \]
% This map is well-defined. It is also clearly an equivalence. Its inverse is given by gluing maps defined on consecutive intervals together to get a map $\mathbb{R} \rightarrow \mathbb{R}\times \Conf(n,I^{N})$. 

% We would like to show that we can forget the two outer terms in the pullback above without changing its homotopy type.
% The source (resp. target) maps 
% \begin{align*}
% p_0: \Phi((-\infty,0]\times \mathbb{R})&\rightarrow \UConf(n,I^N)\times\mathbb{R}\\
% p_p: \Phi([0,\infty)\times \mathbb{R})&\rightarrow \UConf(n,I^N)\times\mathbb{R}
% \end{align*}
% are fibrations, so the pullbacks in the expression above are actually homotopy pullbacks.
% Let us consider the map $q_0:\UConf(n,I^N)\times\mathbb{R}\rightarrow \Phi((-\infty,0]\times I^N)$ which is the identity on $\mathbb{R}$ and sends a configuration to the constant map.
% Clearly, $p_0\circ q_0 = \id$. The following is a homotopy between $q_0\circ p_0$ and $\id$:
% \begin{align*}
% H: \Phi((-\infty,0]\times I^N)\times I &\rightarrow \Phi_0((-\infty,0]\times I^N)\\
% (\phi,t_0,s) &\mapsto (\gamma_s,t_0)
% \end{align*}
% where, for $s<1$:
% $$\gamma_s:t \mapsto \begin{cases}
% &\phi(t-\frac{s}{1-s}),\text{ if }t\in(-\infty,-\frac{s}{1-s}];\\ 
% &\phi(0),\text{ if }t\in[-\frac{s}{1-s},0].
% \end{cases}$$ 
% Notice that this homotopy is only well-defined for $s\in[0,1)$; however, we can set $\gamma_1$ to be the constant path on $\phi(0)$, and this makes the homotopy well-defined and continuous.
% Therefore, the map $p_0$ is an equivalence.
% Similarly, the map $p_p$ is an equivalence.
% When $N=\infty$, the terms $\UConf(n,I^\infty)\times\mathbb{R}$ are contractible, so the projection map 
% $$\Phi((-\infty,0]\times I^N)\times_{\UConf(n,I^\infty)\times\mathbb{R}}\times N_p\C_\infty\times_{\mathbb{R}\times \UConf(n,I^\infty)} \Phi_0([0,\infty) \times I^N)\rightarrow N_p\C_\infty$$
% is a weak equivalence.
% Therefore, there is a weak equivalence $\Phi_0^\infty \simeq N_p\C_\infty$.
% Therefore, we have a weak equivalence between $\Phi_0^\infty$ and $B\C_\infty$.
\end{proof}

\end{document}
